\documentclass{article}
\usepackage{amsmath,amssymb,amsthm}
\usepackage{algorithm}
\usepackage[noend]{algpseudocode}
\usepackage{enumitem}
\usepackage{hyperref}
\newtheorem{theorem}{Theorem}
\newtheorem{lemma}{Lemma}
\newtheorem{assumption}{Assumption}
\newcommand{\prox}{\operatorname{prox}}
\newcommand{\R}{\mathbb{R}}

\begin{document}

\section*{Proximal Gradient Method (PGM)}

\section*{Mathematical Formulation}
Let 
\[
F(x)=f(x)+g(x),\qquad f:\R^{d}\to\R\ \text{smooth},\ g:\R^{d}\to\R\cup\{+\infty\}\ \text{convex (possibly non‑smooth)}.
\]
Denote by $\nabla f$ the gradient of $f$ and assume it is $L$‑Lipschitz, i.e.
\[
\|\nabla f(x)-\nabla f(y)\|\le L\|x-y\|,\qquad\forall x,y\in\R^{d}.
\]
For a stepsize $\eta>0$, the \emph{proximal gradient update} is
\begin{equation}
\boxed{%
x_{k+1}= \prox_{\eta g}\!\bigl(x_{k}-\eta\nabla f(x_{k})\bigr)},
\qquad k=0,1,\dots
\tag{1}
\end{equation}
where the proximal operator of $g$ with parameter $\eta$ is
\[
\prox_{\eta g}(v)=\arg\min_{u\in\R^{d}}\Bigl\{ g(u)+\frac{1}{2\eta}\|u-v\|^{2}\Bigr\}.
\]

\section*{Pseudocode}
\begin{algorithm}[H]
\caption{Proximal Gradient Method}
\begin{algorithmic}[1]
\Require Initial point $x_{0}\in\R^{d}$, stepsize $\eta\in(0,1/L]$, max iterations $K$
\For{$k=0$ \textbf{to} $K-1$}
    \State Compute gradient $g_{k}\gets\nabla f(x_{k})$
    \State Take a gradient step $y_{k}\gets x_{k}-\eta g_{k}$
    \State Apply proximal map $x_{k+1}\gets\prox_{\eta g}(y_{k})$
\EndFor
\State \Return $x_{K}$
\end{algorithmic}
\end{algorithm}

\section*{Dry Run Example}
Consider the scalar problem
\[
f(w)=(w-3)^{2},\qquad g\equiv0,\qquad \eta=0.1,\qquad w_{0}=0.
\]
Since $g\equiv0$, $\prox_{\eta g}(v)=v$. The gradient is $\nabla f(w)=2(w-3)$.

\begin{center}
\begin{tabular}{c|c|c|c}
Iteration $k$ & $w_{k}$ & $\nabla f(w_{k})$ & $w_{k+1}=w_{k}-\eta\nabla f(w_{k})$\\ \hline
0 & $0$ & $-6$ & $0-0.1(-6)=0.6$\\
1 & $0.6$ & $2(0.6-3)=-4.8$ & $0.6-0.1(-4.8)=1.08$\\
2 & $1.08$ & $2(1.08-3)=-3.84$ & $1.08-0.1(-3.84)=1.464$\\
\end{tabular}
\end{center}
The corresponding objective values $F(w_{k})=f(w_{k})$ are
\[
F(w_{0})=9,\quad F(w_{1})=(0.6-3)^{2}=5.76,\quad 
F(w_{2})=(1.08-3)^{2}=3.6864,\quad 
F(w_{3})=(1.464-3)^{2}=2.3665.
\]
The sequence clearly descends toward the optimum $w^{\star}=3$.

\newpage

\section*{Assumptions}
\begin{assumption}[Smoothness of $f$]\label{ass:smooth}
$f$ is convex and differentiable with $L$‑Lipschitz gradient:
\[
\|\nabla f(x)-\nabla f(y)\|\le L\|x-y\|,\qquad\forall x,y\in\R^{d}.
\]
\end{assumption}
\begin{assumption}[Convexity of $g$]\label{ass:convex}
$g$ is proper, closed, convex (possibly non‑smooth). Its proximal operator $\prox_{\eta g}$ is well defined for any $\eta>0$.
\end{assumption}
\begin{assumption}[Stepsize]\label{ass:stepsize}
The stepsize satisfies $0<\eta\le 1/L$.
\end{assumption}

\section*{Main Convergence Theorem}
\begin{theorem}[Ergodic $O(1/k)$ rate]\label{thm:conv}
Let $\{x_{k}\}$ be generated by \eqref{1} under Assumptions~\ref{ass:smooth}–\ref{ass:stepsize}. Then for any minimizer $x^{\star}\in\arg\min F$,
\[
F(x_{k})-F(x^{\star})\le\frac{\|x_{0}-x^{\star}\|^{2}}{2\eta\,k},
\qquad\forall k\ge1.
\tag{2}
\]
Equivalently, the averaged iterate $\bar{x}_{k}:=\frac1k\sum_{t=0}^{k-1}x_{t}$ satisfies the same bound.
\end{theorem}

\section*{Theoretical Properties}
\begin{itemize}[leftmargin=*]
\item \textbf{Stability.} The update is a composition of a non‑expansive map (prox) and a gradient step that is a contraction when $\eta\le 1/L$. Hence the whole iteration is non‑expansive, guaranteeing boundedness of $\{x_{k}\}$.
\item \textbf{Hyperparameter Sensitivity.} The stepsize must respect $\eta\le 1/L$; any larger value may violate the descent inequality (see proof). Within this interval, a larger $\eta$ yields a tighter constant in \eqref{2} because the denominator $2\eta$ grows.
\item \textbf{Comparison to Baselines.}
    \begin{enumerate}
        \item With $g\equiv0$, the method reduces to gradient descent and the same $O(1/k)$ rate is recovered.
        \item Compared with subgradient methods for non‑smooth $g$, the proximal step exploits the structure of $g$ and often yields a much smaller constant and faster practical convergence.
    \end{enumerate}
\end{itemize}

\section*{Discussion}
The theorem shows that the proximal gradient method attains the optimal $O(1/k)$ rate for minimizing the sum of a smooth convex function and a convex (possibly non‑smooth) regularizer. The proof hinges on two elementary but powerful tools:
\begin{enumerate}
\item \emph{Descent Lemma} for $L$‑smooth $f$.
\item \emph{Non‑expansiveness} (firm non‑expansiveness) of the proximal operator.
\end{enumerate}
Both hold under the minimal convexity assumptions stated above, making the result widely applicable (e.g., Lasso, total variation denoising).

\newpage

\section*{Preliminaries (Definitions and Lemmas)}
\begin{definition}[Proximal operator]
For a convex $g$ and $\eta>0$,
\[
\prox_{\eta g}(v)=\arg\min_{u\in\R^{d}}\Bigl\{g(u)+\frac{1}{2\eta}\|u-v\|^{2}\Bigr\}.
\]
\end{definition}
\begin{lemma}[Firm non‑expansiveness]\label{lem:firm}
For any $u,v\in\R^{d}$,
\[
\|\prox_{\eta g}(u)-\prox_{\eta g}(v)\|^{2}
\le \langle u-v,\;\prox_{\eta g}(u)-\prox_{\eta g}(v)\rangle .
\tag{3}
\]
\end{lemma}
\begin{proof}
Standard result from convex analysis; follows from the optimality condition of the proximal subproblem. ∎
\end{proof}
\begin{lemma}[Descent Lemma for $L$‑smooth $f$]\label{lem:descent}
For any $x,y\in\R^{d}$,
\[
f(y)\le f(x)+\langle\nabla f(x),y-x\rangle+\frac{L}{2}\|y-x\|^{2}.
\tag{4}
\]
\end{lemma}
\begin{proof}
Direct consequence of the $L$‑Lipschitz gradient property. ∎
\end{proof}

\section*{Main Proof (Step‑by‑Step Derivation)}
We prove Theorem~\ref{thm:conv}. Throughout, denote $x^{\star}\in\arg\min F$.

\noindent\textbf{Step 1.} Write the optimality condition of the proximal map at iteration $k$.
\[
x_{k+1}= \prox_{\eta g}\!\bigl(x_{k}-\eta\nabla f(x_{k})\bigr)
\Longleftrightarrow
\frac{1}{\eta}\bigl(x_{k}-\eta\nabla f(x_{k})-x_{k+1}\bigr)\in\partial g(x_{k+1}).
\tag{5}
\]

\noindent\textbf{Step 2.} Use convexity of $g$ and (5) to bound $g(x_{k+1})-g(x^{\star})$.
\[
g(x_{k+1})-g(x^{\star})
\le \Big\langle\frac{1}{\eta}\bigl(x_{k}-\eta\nabla f(x_{k})-x_{k+1}\bigr),\;x_{k+1}-x^{\star}\Big\rangle .
\tag{6}
\]

\noindent\textbf{Step 3.} Expand the inner product in (6):
\[
\begin{aligned}
\text{RHS of (6)} 
&= \frac{1}{\eta}\langle x_{k}-x_{k+1},x_{k+1}-x^{\star}\rangle
      -\langle \nabla f(x_{k}),x_{k+1}-x^{\star}\rangle .
\end{aligned}
\tag{7}
\]

\noindent\textbf{Step 4.} Apply the identity $2\langle a,b\rangle =\|a\|^{2}+\|b\|^{2}-\|a-b\|^{2}$ with
$a=x_{k}-x^{\star}$ and $b=x_{k+1}-x^{\star}$ to the first term of (7):
\[
\frac{1}{\eta}\langle x_{k}-x_{k+1},x_{k+1}-x^{\star}\rangle
= \frac{1}{2\eta}\bigl(\|x_{k}-x^{\star}\|^{2}
                     -\|x_{k+1}-x^{\star}\|^{2}
                     -\|x_{k+1}-x_{k}\|^{2}\bigr).
\tag{8}
\]

\noindent\textbf{Step 5.} Combine (6)–(8) and add $f(x_{k+1})-f(x^{\star})$ to both sides:
\[
\begin{aligned}
F(x_{k+1})-F(x^{\star})
&\le f(x_{k+1})-f(x^{\star})
   -\langle \nabla f(x_{k}),x_{k+1}-x^{\star}\rangle \\
&\quad +\frac{1}{2\eta}\bigl(\|x_{k}-x^{\star}\|^{2}
                     -\|x_{k+1}-x^{\star}\|^{2}
                     -\|x_{k+1}-x_{k}\|^{2}\bigr).
\end{aligned}
\tag{9}
\]

\noindent\textbf{Step 6.} Use the Descent Lemma (Lemma~\ref{lem:descent}) with $x=x_{k}$ and $y=x_{k+1}$:
\[
f(x_{k+1})\le f(x_{k})+\langle\nabla f(x_{k}),x_{k+1}-x_{k}\rangle
            +\frac{L}{2}\|x_{k+1}-x_{k}\|^{2}.
\tag{10}
\]

\noindent\textbf{Step 7.} Substitute (10) into (9) and cancel the term $\langle\nabla f(x_{k}),x_{k+1}-x_{k}\rangle$:
\[
\begin{aligned}
F(x_{k+1})-F(x^{\star})
&\le f(x_{k})-f(x^{\star})
   -\langle \nabla f(x_{k}),x_{k}-x^{\star}\rangle \\
&\quad +\frac{1}{2\eta}\bigl(\|x_{k}-x^{\star}\|^{2}
                     -\|x_{k+1}-x^{\star}\|^{2}
                     -\|x_{k+1}-x_{k}\|^{2}\bigr) \\
&\quad +\frac{L}{2}\|x_{k+1}-x_{k}\|^{2}.
\end{aligned}
\tag{11}
\]

\noindent\textbf{Step 8.} Apply convexity of $f$:
\[
f(x_{k})-f(x^{\star})\le\langle \nabla f(x_{k}),x_{k}-x^{\star}\rangle .
\tag{12}
\]
Insert (12) into (11); the first two terms cancel, leaving
\[
F(x_{k+1})-F(x^{\star})
\le \frac{1}{2\eta}\bigl(\|x_{k}-x^{\star}\|^{2}
                     -\|x_{k+1}-x^{\star}\|^{2}
                     -\|x_{k+1}-x_{k}\|^{2}\bigr)
   +\frac{L}{2}\|x_{k+1}-x_{k}\|^{2}.
\tag{13}
\]

\noindent\textbf{Step 9.} Group the norm terms:
\[
F(x_{k+1})-F(x^{\star})
\le \frac{1}{2\eta}\bigl(\|x_{k}-x^{\star}\|^{2}
                     -\|x_{k+1}-x^{\star}\|^{2}\bigr)
   +\Bigl(\frac{L}{2}-\frac{1}{2\eta}\Bigr)\|x_{k+1}-x_{k}\|^{2}.
\tag{14}
\]

\noindent\textbf{Step 10.} By Assumption~\ref{ass:stepsize}, $\eta\le 1/L$, hence $L/2-1/(2\eta)\le0$. Therefore the second term in (14) is non‑positive and can be dropped:
\[
F(x_{k+1})-F(x^{\star})
\le \frac{1}{2\eta}\bigl(\|x_{k}-x^{\star}\|^{2}
                     -\|x_{k+1}-x^{\star}\|^{2}\bigr).
\tag{15}
\]

\noindent\textbf{Step 11.} Rearranging (15) yields the fundamental descent inequality
\[
\|x_{k+1}-x^{\star}\|^{2}
\le \|x_{k}-x^{\star}\|^{2}
      -2\eta\bigl(F(x_{k+1})-F(x^{\star})\bigr).
\tag{16}
\]

\noindent\textbf{Step 12.} Sum (16) for $k=0,\dots, K-1$ (telescoping):
\[
\|x_{K}-x^{\star}\|^{2}
\le \|x_{0}-x^{\star}\|^{2}
      -2\eta\sum_{k=0}^{K-1}\bigl(F(x_{k+1})-F(x^{\star})\bigr).
\tag{17}
\]

\noindent\textbf{Step 13.} Since the left‑hand side is non‑negative,
\[
2\eta\sum_{k=1}^{K}\bigl(F(x_{k})-F(x^{\star})\bigr)
\le \|x_{0}-x^{\star}\|^{2}.
\tag{18}
\]

\noindent\textbf{Step 14.} Divide (18) by $2\eta K$ and use the fact that the minimum of a set is bounded by its average:
\[
\min_{1\le k\le K}\bigl(F(x_{k})-F(x^{\star})\bigr)
\le \frac{1}{K}\sum_{k=1}^{K}\bigl(F(x_{k})-F(x^{\star})\bigr)
\le \frac{\|x_{0}-x^{\star}\|^{2}}{2\eta K}.
\tag{19}
\]

\noindent\textbf{Step 15.} Because $F(x_{k})$ is non‑increasing (from (15)), the same bound holds for every $k\ge1$:
\[
F(x_{k})-F(x^{\star})\le\frac{\|x_{0}-x^{\star}\|^{2}}{2\eta k},
\qquad\forall k\ge1,
\]
which is exactly \eqref{2}. ∎

\section*{Rate Derivation (With Constants)}
The bound \eqref{2} can be written as
\[
F(x_{k})-F^{\star}\le \frac{C}{k},
\qquad C:=\frac{\|x_{0}-x^{\star}\|^{2}}{2\eta},
\]
explicitly displaying the dependence on the stepsize $\eta$ and the initial distance to the optimum.

\section*{Special Cases}
\begin{enumerate}
\item \textbf{Pure Gradient Descent ($g\equiv0$).} The proximal step reduces to the identity, and the proof collapses to the classic $O(1/k)$ result for smooth convex minimization.
\item \textbf{Indicator Regularizer.} If $g=\iota_{\mathcal{C}}$ for a closed convex set $\mathcal{C}$, then $\prox_{\eta g}$ is the Euclidean projection onto $\mathcal{C}$. The method becomes \emph{projected gradient descent}, and the same rate holds.
\item \textbf{Strongly Convex $F$.} If $F$ is $\mu$‑strongly convex, a refined analysis (omitted for brevity) yields a linear rate $O\bigl((1-\eta\mu)^{k}\bigr)$.
\end{enumerate}

\end{document}